\begin{explanationbox}
\textbf{\textcolor{CExplanation}{一句话直觉}}

\textbf{诱导公式是单位圆对称性的代数表达；公式对任意角 $\alpha$ 成立，而口诀“奇变偶不变，符号看象限”是快速判断函数名与符号的做题方法。}

\medskip
\textbf{\textcolor{CExplanation}{核心拆解}}
\begin{description}[style=nextline, leftmargin=2.8em, labelsep=0.5em, itemsep=0.45em]
    \item[\textbf{要点一（奇偶判定）：}] 将角写成 $\frac{k\pi}{2}\pm\alpha$。$k$ 偶则函数名不变；$k$ 奇则正弦与余弦互换、正切与余切互换。
    \item[\textbf{要点二（符号判定）：}] 使用口诀判号时，把 $\alpha$ 暂看作锐角或参考角，再看整体角 $\frac{k\pi}{2}\pm\alpha$ 所在象限，取原函数在该象限的符号。
    \item[\textbf{要点三（任意角与参考角）：}] 公式中的 $\alpha$ 可以是任意角；“把 $\alpha$ 看作锐角”只服务于口诀判号，不能理解为公式本身只对锐角成立。
    \item[\textbf{要点四（定义域）：}] 含 $\tan$ 或 $\cot$ 的公式必须检查左右两边是否有意义。出现 $\tan\theta$ 要求 $\cos\theta\neq0$，出现 $\cot\theta$ 要求 $\sin\theta\neq0$。
\end{description}

\medskip
\textbf{\textcolor{CExplanation}{几何本质（对称变换）}}

角 $\alpha$ 对应单位圆上点 $P(\cos\alpha,\sin\alpha)$。$-\alpha$ 关于 $x$ 轴对称，对应点 $(\cos\alpha,-\sin\alpha)$；$\pi-\alpha$ 关于 $y$ 轴对称，对应点 $(-\cos\alpha,\sin\alpha)$；$\pi+\alpha$ 关于原点对称，对应点 $(-\cos\alpha,-\sin\alpha)$。

对于变名公式，$\frac{\pi}{2}-\alpha$ 对应点 $(\sin\alpha,\cos\alpha)$，可理解为关于直线 $y=x$ 对称；$\frac{\pi}{2}+\alpha$ 对应点 $(-\sin\alpha,\cos\alpha)$，可理解为将点 $P$ 逆时针旋转 $90^\circ$。坐标变化直接给出新的正弦、余弦值。

\medskip
\textbf{\textcolor{CExplanation}{代数意义（周期性与对称性）}}

$\sin$、$\cos$ 的周期是 $2\pi$，所以 $2k\pi+\alpha$ 对应的正弦、余弦值不变；$\tan$、$\cot$ 的周期是 $\pi$，所以 $k\pi+\alpha$ 对应的正切、余切值不变。

不能把 $\pi\pm\alpha$ 简单理解为“半周期都变号”。更准确地说，$\pi+\alpha$ 对应绕原点旋转 $180^\circ$，所以 $\sin$、$\cos$ 变号，而 $\tan$、$\cot$ 不变；$\pi-\alpha$ 对应关于 $y$ 轴对称，所以 $\sin$ 不变，$\cos$ 变号，$\tan$、$\cot$ 变号。

\medskip
\textbf{\textcolor{CExplanation}{考点价值（化简工具）}}

诱导公式可以把任意角三角函数化为同名或异名的参考角三角函数；在求特殊角数值时，再进一步化为锐角三角函数。
\begin{itemize}[leftmargin=*, itemsep=0.5em, topsep=0.4em]
    \item \textbf{考法一（直接求值）：} 给定角（如 $\frac{5\pi}{4}$），先化为参考角，再利用特殊角三角函数值计算。
    \item \textbf{考法二（恒等式化简）：} 通过多次使用诱导公式，把复杂角统一成同一类角，便于合并、约分或证明。
\end{itemize}

\medskip
\textbf{\textcolor{CExplanation}{顿悟点}}

\textbf{公式本身允许 $\alpha$ 为任意角；口诀判号时才把 $\alpha$ 暂看作锐角或参考角。这样既不丢严谨性，又能快速判断符号。}

\medskip
\textbf{\textcolor{CExplanation}{使用场景}}
\begin{itemize}[leftmargin=*, itemsep=0.5em, topsep=0.4em]
    \item \textbf{场景一（化简求值）：} 任意角化归为参考角或锐角。
    \item \textbf{场景二（三角变形）：} 在恒等式变换中调整角度形式。
    \item \textbf{场景三（三角方程）：} 化简含 $\tan$、$\cot$ 的式子时，要先检查定义域，避免增根。
\end{itemize}
\end{explanationbox}
